\appendix

\section{Claim-to-Public-Artifact Matrix}
\label{app:public-configs}

\begin{table}[H]
  \centering
  \caption{Public files associated with the report's central claims.}
  \begin{tabularx}{\textwidth}{@{}L{0.22\textwidth}L{0.30\textwidth}Y@{}}
    \toprule
    Purpose & Public artifact & Reproducibility boundary \\
    \midrule
    Portable numerical correctness &
      \path{tests/test_portable_backend_gpu.py} and
      \path{artifact/configs/smoke-synthetic.toml} &
      Compares FP32 and FP16 cache materialization and the full
      $K\Delta\alpha$ update with an explicit RBF oracle; also exercises the
      patched solver end to end. \\
    Matched performance protocol &
      \path{artifact/configs/ember65k-ab.toml} and
      \path{artifact/reference/pro6000.json} &
      Encodes equal logical cache capacity and balanced repetitions. The
      public portable backend does not reproduce the 1.351$\times$
      architecture-tuned result. \\
    Held-out quality protocol &
      \path{artifact/configs/ember65k-quality.toml} &
      Runs stock FP32 and coherent FP16 on one held-out split and summarizes
      accuracy, AUC, disagreements, McNemar's exact test, and score
      correlation. It is not an equivalence proof. \\
    Work--depth analysis &
      \path{studies/work-depth/metrics.py} &
      Separates global columns, local rows, and dependent decision stages.
      The private captured traces used in the report are not redistributed. \\
    \bottomrule
  \end{tabularx}
\end{table}

The source release is intentionally smaller than the development workspace.
Generated datasets, binaries, profiler reports, and failed experimental
branches are not required to understand or exercise the mechanism.

\section{Repository Layout and Quick Start}
\label{app:quick-start}

The fresh public repository has the following stable surface:

\begin{lstlisting}[style=codeblock]
coherent-svm/
|-- src/portable_fp16_rbf_backend.cu
|-- include/
|-- patches/cuml-26.06.00/
|-- coherent_svm/
|-- artifact/configs/
|-- artifact/schemas/
|-- studies/work-depth/
`-- tests/
\end{lstlisting}

The repository is public; only dependency paths remain machine-specific:

\begin{lstlisting}[style=codeblock]
git clone https://github.com/yongqiu-star/coherent-svm
cd coherent-svm
uv sync --frozen --extra gpu --extra quality --extra dev
uv run coherent-svm doctor

export CUML_SOURCE=<path-to-cuml-26.06.00>
export CUVS_INCLUDE=<path-to-cuvs-include>
git -C "$CUML_SOURCE" apply \
  "$PWD/patches/cuml-26.06.00/coherent-fp16.patch"
PYTHON="$PWD/.venv/bin/python" scripts/build.sh

COHERENT_SVM_LIBRARY=build/libcoherent_svm.so \
  uv run pytest -q
uv run coherent-svm reproduce \
  --config artifact/configs/smoke-synthetic.toml \
  --library build/libcoherent_svm.so \
  --output-root results
\end{lstlisting}

The main prepared-data protocol replaces the smoke configuration with
\path{artifact/configs/ember65k-ab.toml} and adds
\code{--data-root <prepared-data-root>}. Dataset arrays remain external to the
repository; their normalized layout is documented in
\path{artifact/README.md}. The release specifies the prepared-array contract
but does not include EMBER's feature-extraction pipeline.
The complete patched-solver reproduction path currently requires compute
capability 12.0; the standalone portable-backend oracle test builds for the
local CUDA architecture.

\section{Normalized Output Contract}
\label{app:output-contract}

Each reproduction writes five public files:

\begin{lstlisting}[style=codeblock]
results/<config-id>/<run-id>/
|-- manifest.json
|-- measurements.jsonl
|-- summary.json
|-- checks.json
`-- artifacts.sha256
\end{lstlisting}

The manifest retains the public source revision, dirty state, configuration
and library digests, logical dataset identifier, Python version, and platform
class. It intentionally omits the hostname, user name, GPU UUID, private Git
remote, and absolute command paths. Shape, device model, compute capability,
solver parameters, timing, iteration count, and sampled training accuracy are
written with each measurement rather than hidden in a prose log.

For example, the balanced performance summary has the following stable
shape; absolute times are observations rather than cross-machine pass
conditions:

\begin{lstlisting}[style=codeblock]
{
  "config_id": "ember65k-ab",
  "baseline": {"method": "direct_fp32",
               "median_fit_seconds": "<measured>"},
  "candidate": {"method": "coherent_fp16",
                "median_fit_seconds": "<measured>"},
  "speedup": "<ratio>",
  "speedup_ci95": ["<low>", "<high>"],
  "repeats": 7
}
\end{lstlisting}

\section{Licensing and Redistribution Boundary}
\label{app:licensing}

The original reference backend, runner, configurations, schemas, tests, and
documentation are released under Apache-2.0. Changes to cuML are distributed
as patches against cuML 26.06.00 and retain the upstream license notices when
applied. CUDA, cuBLAS, cuML, RAPIDS dependencies, and datasets are external;
the repository neither vendors their binaries nor redistributes raw data.

The optimized \code{sm\_120a} backend is excluded because file-level
provenance and redistribution rights for all third-party-derived components
have not yet been resolved. This is a release boundary, not a performance
claim: the portable backend validates the coherent operator and integration,
whereas the 1.351$\times$ measurement describes the unreleased tuned backend.
A later release of that backend requires an independent provenance and
licensing review.
